

\section{A technique to complete the solution to the generalized road coloring conjecture}  
\label{sec:4.6}

Using the result from the previous section, a technique suggests itself to solve the general problem.  
Suppose the graph is periodic of order $p$.  
Look at the collection of all minimal coloring semigroups.  
Suppose the kernels are of rank $t$.  
There are two possibilities.  
If one of the minimal kernels is a right group, from the previous result, we immediately deduce that $p=t$.  
If no minimal coloring semigroup has a kernel which is a right group, one must show that some recoloring of one of the minimal semigroups produces a kernel of rank strictly less than $t$.  

Thus, if one can prove the following conjecture, the proof of the generalized road coloring problem would be complete.

\begin{conjecture} 	\label{con:4.2}
Let $S$ be a coloring semigroup.  
Suppose the kernel of $S$ is not a right group.  
Then $S$ is not a minimal coloring semigroup.
\end{conjecture}

Consider the following example of recoloring in the instance that a kernel is not a right group.

\begin{example}  	\label{exm:4.2}
Let $G$ be given by the following adjacency matrix which is colored
as indicated
\[
	A = \bmatrix
	   0 &1 &0 &1 \\
	   0 &0 &1 &1 \\
	   1 &1 &0 &0 \\
	   1 &0 &1 &0 \endbmatrix
	=
		\bmatrix
		0 &0 &0 &1 \\
		0 &0 &0 &1 \\
		1 &0 &0 &0 \\
		1 &0 &0 &0 \endbmatrix
	+
		\bmatrix
		0 &1 &0 &0 \\
		0 &0 &1 &0 \\
		0 &1 &0 &0 \\
		0 &0 &1 &0 \endbmatrix = r_1 + b_1
\]
Let $S_1 = \langle r_1,b_1 \rangle$.  Then $S_1$ is, in fact, a completely simple
semigroup of matrices of rank two.  $S_1$ consists of the following eight matrices.
\begin{equation*}
	S_1 = \left \{
		\begin{aligned}
		\bmatrix
		 1 &0 &0 &0 \\
		 1 &0 &0 &0 \\
 		 0 &0 &0 &1 \\
		 0 &0 &0 &1 \endbmatrix, 
		\bmatrix
		 0 &0 &0 &1 \\
		 0 &0 &0 &1 \\
		 1 &0 &0 &0 \\
		 1 &0 &0 &0 \endbmatrix, 
		\bmatrix
		 0 &1 &0 &0 \\
		 0 &1 &0 &0 \\
		 0 &0 &1 &0 \\
		 0 &0 &1 &0 \endbmatrix,
		\bmatrix
		 0 &0 &1 &0 \\
		 0 &0 &1 &0 \\
		 0 &1 &0 &0 \\
		 0 &1 &0 &0 \endbmatrix	\\
		\bmatrix
		1 &0 &0 &0 \\
		0 &0 &0 &1 \\
		1 &0 &0 &0 \\
		0 &0 &0 &1 \endbmatrix,
		\bmatrix
		0 &0 &0 &1 \\
		1 &0 &0 &0 \\
		0 &0 &0 &1 \\
		1 &0 &0 &0 \endbmatrix,
		\bmatrix
		0 &0 &1 &0 \\
		0 &1 &0 &0 \\
		0 &0 &1 &0 \\
		0 &1 &0 &0 \endbmatrix,
		\bmatrix
		0 &1 &0 &0 \\
		0 &0 &1 &0 \\
		0 &1 &0 &0 \\
		0 &0 &1 &0 \endbmatrix
		\end{aligned} 
		\right\}
\end{equation*}

	Since $S_1$ (which is its own kernel), is not a right group, the
functions within it do not all have the same partition.  There are two
partitions in this example.
\[
	\pi_1 = \left\{ \{1,2\}, \{3,4\} \right\}, \quad
		\pi_2 = \left\{ \{1,3\}, \{2,4\} \right\}
\]
\end{example}

In general there will be many such partitions.  
There will always be at least one vertex $v$, which is paired with distinct vertices in two of these partitions.  
(In this example, any vertex would suffice.  
That is, 1 is paired with 2 in $\pi_1$, but with 3 in $\pi_2$.)  
Perhaps recoloring at such a vertex $v$ will allow one to reduce the rank.
Computer experiments offer some hope this is true.
Certainly, in this case, recoloring at 1 produces the following.
\[
	\bmatrix
	0 &1 &0 &1 \\
	0 &0 &1 &1 \\
	1 &1 &0 &0 \\
	1 &0 &1 &0 \endbmatrix
	=
	\bmatrix
	0 &1 &0 &0 \\
	0 &0 &0 &1 \\
	1 &0 &0 &0 \\
	1 &0 &0 &0 \endbmatrix
	+
	\bmatrix
	0 &0 &0 &1 \\
	0 &0 &1 &0 \\
	0 &1 &0 &0 \\
	0 &0 &1 &0 \endbmatrix
	= r_2 + b_2
\]
A simple computation will show that
\[
	r_2b_2r_2 = \bmatrix
		1 &0 &0 &0 \\
		1 &0 &0 &0 \\
		1 &0 &0 &0 \\ 
		1 &0 &0 &0 \endbmatrix.
\]
Thus $S_2 = \langle r_2,b_2 \rangle$ contains a rank one matrix, or
equivalently, a synchronizing instruction.
